% ======================================================================
\section{Discussion: What Matters for Lunar VPR}\label{sec:discussion}
% ======================================================================
Viewing both stages as OT problems lets us separate three kinds of design decision and rank them by
their measured effect on LunarSim-VPR.

\paragraph*{Training data first}
Illumination diversity in training is the largest single factor: it is worth $\fixIllDrop$ points of
R@1 with a frozen backbone and $\ftFixIllDrop$ with a fine-tuned one, more than any architectural
choice. Because a more adaptable model fits a single training light more closely, illumination
diversity and backbone fine-tuning have to go together.

\paragraph*{Then backbone adaptation, then aggregation}
Fine-tuning the last four blocks adds $11$--$19$ points for every head. The aggregation matters only
at the coarse level (cluster-based aggregation beats global pooling by more than $20$ points) and,
once the backbone is trained, at the level of residual vs.\ OT aggregation. Within the OT family the
assignment is a weak design lever: all members are equivalent on a frozen backbone (P1, P2), and on
a fine-tuned backbone the end points are preferable to a learned intermediate value. The cheapest
member, $\lambda=0$, needs no Sinkhorn loop and matches SALAD's balanced assignment within seed noise. A
learnable $\lambda$ adds a parameter that the loss barely constrains (P2), so it should be fixed.

\paragraph*{Fusion, not the matcher}
Local evidence on the Moon is abundant but unreliable, because shadow edges dominate keypoint
detections. Its value lies in the queries where the global ranking is uncertain, which is exactly
where the bounded-override property of the fusion rule (\cref{prop:override}) lets it act. With
fusion in place, a training-free OT matcher, plain mutual nearest neighbours and a learned
terrestrial matcher perform alike, so the cheapest adequate scorer, MNN+MAGSAC++, is a sensible default
(LightGlue costs $\approx\tLgK$\,ms and FGW $\approx\tFgwK$\,ms per query on a CPU). The structural
FGW matcher earns its cost mainly where no learned matcher is available or trusted.

\paragraph*{What the theory did and did not predict}
\Cref{prop:colinv} correctly predicted that, with fixed tokens, the assignment family collapses to
one operating point and that a learned $\lambda$ would not move. It does not predict the ordering
once the tokens are trained; that ordering, and the advantage of residual aggregation, are empirical
findings on simulated data that should be re-examined on real imagery.
